% Canonical Paper IV appendix.
The tables below make the finite calibrations reproducible without changing
their interpretation.

\subsection{The finite-field projective tower}

\[
\begin{array}{c@{\qquad}c@{\qquad}c}
|\PP^1(\F_7)|=8,
&|\Biv_{\F_7}^{\mathrm{reg}}|=8\cdot7=56,
&|\PGL_2(\F_7)|=8\cdot7\cdot6=336.
\end{array}
\]
Every bivaluation has six projective-frame lifts.

\subsection{Binary Horner words of length four}

For $k=2$, $n=4$, and Hamming weight two:
\[
\begin{array}{c@{\qquad}c}
\toprule
\text{word }(d_1d_2d_3d_4) & \text{constant in }16x+c\\
\midrule
0011&3\\
0101&5\\
0110&6\\
1001&9\\
1010&10\\
1100&12\\
\bottomrule
\end{array}
\]
All six words have charge $(2,4\log2)$.

For $n=8$, the fixed-weight counts are:
\[
\begin{array}{c|rrrrrrrrr}
m&0&1&2&3&4&5&6&7&8\\
\hline
\binom{8}{m}&1&8&28&56&70&56&28&8&1
\end{array}
\]
The central charge class therefore contains seventy distinct affine operators
in characteristic zero.

\subsection{Free-group ball count}

For the free group $F_r$ with the standard symmetric alphabet of size $2r$,
the sphere of radius $n\ge1$ has
\[
  2r(2r-1)^{n-1}
\]
reduced words.  Hence
\[
  |B(n)|
  =1+2r\frac{(2r-1)^n-1}{2r-2}.
\]
The exponential metric growth coexists with the linear-time stack algorithm
for freely reducing one input word.

\subsection{Finite-field roots and transform networks}

Repeated squaring gives
\[
  3^{128}\equiv256\equiv-1\pmod{257},
  \qquad 3^{256}\equiv1\pmod{257}.
\]
Thus $3$ has order $256$.  The selected roots and network counts are:

\begin{center}
\begin{tabular}{rrrrrr}
\toprule
$N$ & $\zeta_N=3^{256/N}$ & layers & butterflies & $\log_2$ state count & direct bits\\
\midrule
8  &64  &3&12&$8\log_2 257$&72\\
16 &249 &4&32&$16\log_2 257$&144\\
\bottomrule
\end{tabular}
\end{center}

The ``direct bits'' column uses nine bits per field element.  The logarithmic
cardinality is smaller because $257$ is not a power of two.  Twiddle constants
are not included in either dynamic-state column.

\subsection{Matrix-chain calibration}

For dimensions $(100,10,1,2)$:
\[
\begin{array}{c@{\qquad}c@{\qquad}c}
\toprule
plan & \text{scalar multiplications} & \text{principal intermediate entries}\\
\midrule
$(AB)C$ &100\cdot10\cdot1+100\cdot1\cdot2=1200&100\\
$A(BC)$ &10\cdot1\cdot2+100\cdot10\cdot2=2020&20\\
\bottomrule
\end{array}
\]
Peak memory cannot be recovered from the last column without an allocation and
retention policy.

\subsection{Reverse-chain endpoint schedules}

Using Proposition~\ref{prop:reverse-chain-extremes}:
\[
\begin{array}{c@{\qquad}c@{\qquad}c@{\qquad}c}
\toprule
$N$ & \text{store-all work} & \text{naive recompute work} & \text{extra prefix work}\\
\midrule
8  &16&44&28\\
16 &32&152&120\\
\bottomrule
\end{array}
\]
The work convention charges one unit per forward local evaluation and one per
reverse local step.  It excludes instruction, communication, and variable-size
state costs.
